%% toc: P2 23 23.1
\section{Proofs for the further results on Church's simple theory of types}
\label{sec:annex-p02-w02c-p02-028}

Notation and conventions are those of \cref{sec:annex-p02-w02-p02-040}: (P)
is the propositional input proved in
\cref{annex:p02-w02-p02-040-church-prop-complete}, ``(P), V'' means an
instance of a tautology followed by applications of Rule~V, and bound
variables are renamed as described there without further mention.  For a
set \(\Sigma\) of axioms, \(\Gamma\derives_\Sigma\mathbf B_o\) means that
there is a proof of \(\mathbf B_o\) on the assumptions \(\Gamma\) whose
axioms are variants of members of \(\Sigma\).  The proof of the deduction
theorem in \cref{annex:p02-w02-p02-040-church-deduction} uses, besides the
proof it transforms, only theorems obtained from Axioms 1--6\(^\gamma\); so
it holds for \(\derives_\Sigma\) whenever \(\Sigma\) contains Axioms
1--6\(^\gamma\).  The same holds for the lemma
\cref{thm:p02-w02-p02-040-church-12-21} in its clauses 12--20, which uses
Axioms 1--6\(^\gamma\) only.

\begin{annexitem}{annex:p02-w02c-p02-028-predecessor}{thm:p02-w02c-p02-028-predecessor}{No description-free same-type predecessor}
\begin{proof}
\emph{Step 1: two models and a map between their individuals.}  Since
\(D\) has at least two members, fix \(d_0\ne d_1\) in \(D\).  Let
\(E=\{e_0,e_1\}\) with \(e_0\ne e_1\), let \(g:D\to E\) send \(d_0\) to
\(e_0\) and every other member of \(D\) to \(e_1\), and let \(s:E\to D\)
send \(e_0\) to \(d_0\) and \(e_1\) to \(d_1\).  Then \(g\) is onto and
\(g(s(e))=e\) for every \(e\in E\).  Write \(\mathcal D_\gamma\) for the
domains of \(\mathcal S(D,c)\) and \(\mathcal E_\gamma\) for those of
\(\mathcal S(E,c')\), with \(c'\) any family of choice functions for \(E\).
Write \(\sem{\cdot}^D\) and \(\sem{\cdot}^E\) for denotation in the two
models.  The model \(\mathcal S(E,c')\) is the interpretation with two
individuals, in which Axiom~8 is not valid
(\cref{thm:p02-w02-p02-040-church-soundness}, clause~1).

\emph{Step 2: a relation between the models.}  Define relations
\(R_\gamma\subseteq\mathcal D_\gamma\times\mathcal E_\gamma\) by recursion
on types:
\[
 R_\iota=\{(d,g(d)):d\in D\},\qquad
 R_o=\mathcal D_o\times\mathcal E_o,
\]
and \((u,v)\in R_{\alpha\beta}\) iff \((u(a),v(b))\in R_\alpha\) for every
\((a,b)\in R_\beta\).  Each \(R_\gamma\) is nonempty: \(R_\iota\) and \(R_o\)
are, and if \((a_0,b_0)\in R_\alpha\) then the constant functions with
values \(a_0\) and \(b_0\) are related by \(R_{\alpha\beta}\).

\emph{Step 3: types ending in \(o\).}  Say that \(\gamma\) \emph{ends in}
\(o\) if \(\gamma=o\) or \(\gamma=(\alpha\beta)\) with \(\alpha\) ending in
\(o\).  If \(\gamma\) ends in \(o\), then \(R_\gamma\) is all of
\(\mathcal D_\gamma\times\mathcal E_\gamma\).  For \(\gamma=o\) this is the
definition.  If \(\gamma=(\alpha\beta)\) with \(\alpha\) ending in \(o\),
then for any \(u,v\) and any \((a,b)\in R_\beta\) the pair \((u(a),v(b))\)
lies in \(R_\alpha\), which is everything by the induction hypothesis; so
\((u,v)\in R_{\alpha\beta}\).  The types \(oo\), \(ooo\) and \(o(o\beta)\)
of \(N_{oo}\), \(A_{ooo}\) and \(\Pi_{o(o\beta)}\) end in \(o\).

\emph{Step 4: description-free wffs respect \(R\).}  Call assignments
\(t\) into \(\mathcal S(D,c)\) and \(t'\) into \(\mathcal S(E,c')\)
\emph{related} if \((t(\mathbf x_\gamma),t'(\mathbf x_\gamma))\in R_\gamma\)
for every variable \(\mathbf x_\gamma\); related pairs exist by Step~2.  We
show by induction on the formation of a wff \(\mathbf M_\gamma\) in which no
selection constant occurs that
\((\sem{\mathbf M_\gamma}^D_t,\sem{\mathbf M_\gamma}^E_{t'})\in R_\gamma\)
for all related \(t,t'\).
\begin{itemize}
 \item A variable: this is the hypothesis on \(t,t'\).
 \item A constant: the only constants in \(\mathbf M_\gamma\) are
   \(N_{oo}\), \(A_{ooo}\) and \(\Pi_{o(o\beta)}\), and their types end in
   \(o\), so Step~3 applies.
 \item An application \((\mathbf F_{\alpha\beta}\mathbf A_\beta)\): by the
   induction hypothesis the denotations of \(\mathbf F_{\alpha\beta}\) are
   related by \(R_{\alpha\beta}\) and those of \(\mathbf A_\beta\) by
   \(R_\beta\), so by the definition of \(R_{\alpha\beta}\) the values are
   related by \(R_\alpha\).
 \item An abstraction \((\lambda\mathbf x_\beta\mathbf M_\alpha)\): let
   \((a,b)\in R_\beta\).  The assignments \(t[\mathbf x_\beta\mapsto a]\) and
   \(t'[\mathbf x_\beta\mapsto b]\) are related, so by the induction
   hypothesis \(\sem{\mathbf M_\alpha}^D_{t[\mathbf x_\beta\mapsto a]}\) and
   \(\sem{\mathbf M_\alpha}^E_{t'[\mathbf x_\beta\mapsto b]}\) are related by
   \(R_\alpha\).  These are the values of the two denotations of the
   abstract at \(a\) and at \(b\), so the denotations are related by
   \(R_{\alpha\beta}\).
\end{itemize}
This is the only step that uses the hypothesis on \(\mathbf F_{\iota'\iota'}\).
A selection constant \(\iota_{\beta(o\beta)}\) would have to be treated as a
constant, and the choice functions \(c_\beta\) and \(c'_\beta\) need not be
related by \(R_{\beta(o\beta)}\), because \(\beta(o\beta)\) does not end in
\(o\).

\emph{Step 5: \(R_{\iota'}\) is single-valued.}  Recall
\(\iota'=((\iota\iota)(\iota\iota))\).  Suppose \((u,v)\in R_{\iota'}\) and
\((u,w)\in R_{\iota'}\); we show \(v=w\).  Let \(h\in\mathcal E_{\iota\iota}\)
and put \(f=s\circ h\circ g\in\mathcal D_{\iota\iota}\).  Then
\((f,h)\in R_{\iota\iota}\): if \((d,g(d))\in R_\iota\), then
\(g(f(d))=g(s(h(g(d))))=h(g(d))\), so \((f(d),h(g(d)))\in R_\iota\).
Hence \((u(f),v(h))\in R_{\iota\iota}\) and \((u(f),w(h))\in R_{\iota\iota}\).
For \(e\in E\) we have \((s(e),e)\in R_\iota\) because \(g(s(e))=e\), so
\[
 (u(f)(s(e)),\,v(h)(e))\in R_\iota\quad\text{and}\quad(u(f)(s(e)),\,w(h)(e))\in R_\iota.
\]
Since \(R_\iota\) is the graph of \(g\), both \(v(h)(e)\) and \(w(h)(e)\)
equal \(g(u(f)(s(e)))\).  So \(v(h)=w(h)\) for every \(h\), and \(v=w\),
since members of \(\mathcal E_{\iota'}\) are functions.

\emph{Step 6: numerals in the two models.}  In any standard model with
domain \(X\) of individuals, unfolding the abbreviations of
\cref{def:p02-w02-p02-040-church-abbrev} and the denotation clauses gives
\(\sem{n_{\iota'}}=(h\mapsto h^n)\), where \(n_{\iota'}\) is the \(n\)-th
numeral and \(h^n\) the \(n\)-fold composite of \(h:X\to X\) (with \(h^0\)
the identity), and
\(\sem{S_{\iota'\iota'}}(u)=(h\mapsto h\circ u(h))\).  Hence
\(\sem{S_{\iota'\iota'}n_{\iota'}}=\sem{(n+1)_{\iota'}}\).  Now take
\(X=E\).  The four functions \(E\to E\) are the identity, the
transposition \(\sigma\) of \(e_0,e_1\), and the constant maps \(\kappa_0\),
\(\kappa_1\) with values \(e_0\), \(e_1\).  Each satisfies \(h^3=h\):
\(\sigma^2\) is the identity, and constant maps and the identity are
idempotent.  Therefore
\[
 \sem{S_{\iota'\iota'}0_{\iota'}}^E=\sem{1_{\iota'}}^E=(h\mapsto h)
 =(h\mapsto h^3)=\sem{3_{\iota'}}^E=\sem{S_{\iota'\iota'}2_{\iota'}}^E.
\]
On the other hand \(\sem{0_{\iota'}}^E(\kappa_0)\) is the identity of \(E\)
and \(\sem{2_{\iota'}}^E(\kappa_0)=\kappa_0\), which differ at \(e_1\); so
\(\sem{0_{\iota'}}^E\ne\sem{2_{\iota'}}^E\).  Thus \(0_{\iota'}\) and
\(2_{\iota'}\) are distinct numbers with the same successor: this is the
failure of Axiom~8 with two individuals.

\emph{Step 7: the contradiction.}  Suppose that, in \(\mathcal S(D,c)\),
\(\sem{\mathbf F_{\iota'\iota'}(S_{\iota'\iota'}n_{\iota'})}^D=\sem{n_{\iota'}}^D\)
for \(n=0\) and \(n=2\).  The wffs \(\mathbf F_{\iota'\iota'}\),
\(S_{\iota'\iota'}\), \(0_{\iota'}\) and \(2_{\iota'}\) are closed and contain
no selection constant, so by Step~4, applied to the closed
description-free wffs \(\mathbf F_{\iota'\iota'}(S_{\iota'\iota'}n_{\iota'})\)
and \(n_{\iota'}\),
\[
 \bigl(\sem{\mathbf F_{\iota'\iota'}(S_{\iota'\iota'}n_{\iota'})}^D,\
 \sem{\mathbf F_{\iota'\iota'}}^E(\sem{S_{\iota'\iota'}n_{\iota'}}^E)\bigr)\in R_{\iota'},
\]
and the first component is \(\sem{n_{\iota'}}^D\) by the supposition.
Also \((\sem{n_{\iota'}}^D,\sem{n_{\iota'}}^E)\in R_{\iota'}\).  By
Step~5,
\(\sem{\mathbf F_{\iota'\iota'}}^E(\sem{S_{\iota'\iota'}n_{\iota'}}^E)=\sem{n_{\iota'}}^E\)
for \(n=0,2\).  By Step~6 the two left-hand sides are equal, so
\(\sem{0_{\iota'}}^E=\sem{2_{\iota'}}^E\), contradicting Step~6.  This proves
the displayed disjunction.

\emph{Step 8: the consequences.}  The formulas
\(N_{o\iota'}0_{\iota'}\) and \(N_{o\iota'}2_{\iota'}\) are theorems of
Axioms 1--6\(^\gamma\), by 22\(^\iota\) and two uses of 23\(^\iota\) with
conversion (\cref{thm:p02-w02-p02-040-church-induction}).  By clause~1 of
\cref{thm:p02-w02-p02-040-church-soundness}, Axioms 1--6\(^\gamma\) are valid
in \(\mathcal S(D,c)\) and the rules preserve validity, so both formulas
are valid there.  If
\(N_{o\iota'}x_{\iota'}\supset\mathbf F_{\iota'\iota'}(S_{\iota'\iota'}x_{\iota'})=x_{\iota'}\)
were valid, then assigning to \(x_{\iota'}\) the denotation of
\(0_{\iota'}\), and then that of \(2_{\iota'}\), would make the consequent
true; by clause~1 of the same theorem, \(=\) denotes identity, so both
equations of Step~7 would hold, which is impossible.  Finally, every
theorem of the system is valid in \(\mathcal S(\N,c)\) (clause~2 of the
same theorem), and \(\N\) has at least two members; so the formula is not a
theorem.
\end{proof}
\end{annexitem}

\begin{annexitem}{annex:p02-w02c-p02-028-classes}{thm:p02-w02c-p02-028-classes}{Coextensive propositional functions determine equal classes}
\begin{proof}
Fix \(\mathbf F_{o\alpha}\), and write \(\mathbf F\) for
\(\mathbf F_{o\alpha}x_\alpha\).  Let \(\mathbf B(\mathbf y)\) be
\((\mathbf F)[\mathbf y=0_{\iota'}]\vee(\sim\mathbf F)[\mathbf y=1_{\iota'}]\)
for a wff \(\mathbf y\) of type \(\iota'\), and let
\(\mathbf t_{\mathbf F}\) be
\(\iota_{\iota'(o\iota')}(\lambda y_{\iota'}\mathbf B(y_{\iota'}))\).  By
Rule~II, \(\mathbf C_{\mathbf F}x_\alpha\) converts to
\(\mathbf t_{\mathbf F}\), and \((\lambda y_{\iota'}\mathbf B(y_{\iota'}))\mathbf y\)
converts to \(\mathbf B(\mathbf y)\).  Rules II and III act on parts of
formulas, so in any formula a part \(\mathbf C_{\mathbf F}x_\alpha\) may be
replaced by \(\mathbf t_{\mathbf F}\) and conversely.  All axioms below are
variants of 1--6\(^\gamma\), 9\(^{\iota'}\), 10\(^{\iota'\alpha}\) and, in
step~(d), 7.

\emph{(a) \(\mathbf F\derives\mathbf C_{\mathbf F}x_\alpha=0_{\iota'}\).}
Take the variant
\(f_{o\iota'}u_{\iota'}\supset.\,(v_{\iota'})[f_{o\iota'}v_{\iota'}\supset u_{\iota'}=v_{\iota'}]
\supset f_{o\iota'}(\iota_{\iota'(o\iota')}f_{o\iota'})\) of 9\(^{\iota'}\),
with \(f_{o\iota'},u_{\iota'},v_{\iota'}\) not occurring in
\(\mathbf F_{o\alpha}\).  By IV\('\) substitute
\(\lambda y_{\iota'}\mathbf B(y_{\iota'})\) for \(f_{o\iota'}\) and
\(0_{\iota'}\) for \(u_{\iota'}\), and convert by Rule~II:
\[
 \derives\mathbf B(0_{\iota'})\supset.\,(v_{\iota'})[\mathbf B(v_{\iota'})\supset0_{\iota'}=v_{\iota'}]
 \supset\mathbf B(\mathbf t_{\mathbf F}).\tag{\(*\)}
\]
First, \(\derives0_{\iota'}=0_{\iota'}\) by 16\(^{\iota'}\), and the
tautology \(q\supset.\,p\supset pq\vee r\) with (P), V gives
\(\mathbf F\derives\mathbf B(0_{\iota'})\).  Second, the tautology
\(p\supset.\,pq\vee(\sim p)r\supset q\) gives
\(\mathbf F\derives\mathbf B(v_{\iota'})\supset v_{\iota'}=0_{\iota'}\);
with 19\(^{\iota'}\) and (P), V,
\(\mathbf F\derives\mathbf B(v_{\iota'})\supset0_{\iota'}=v_{\iota'}\); and
VI\('\) on \(v_{\iota'}\), which is not free in the assumption, gives
\(\mathbf F\derives(v_{\iota'})[\mathbf B(v_{\iota'})\supset0_{\iota'}=v_{\iota'}]\).
Two applications of Rule~V to \((*)\) give
\(\mathbf F\derives\mathbf B(\mathbf t_{\mathbf F})\), and the same
tautology \(p\supset.\,pq\vee(\sim p)r\supset q\) gives
\(\mathbf F\derives\mathbf t_{\mathbf F}=0_{\iota'}\).  Replacing
\(\mathbf t_{\mathbf F}\) by \(\mathbf C_{\mathbf F}x_\alpha\) (Rule~III)
gives (a).

\emph{(b) \(\sim\mathbf F\derives\mathbf C_{\mathbf F}x_\alpha=1_{\iota'}\).}
The same argument with \(1_{\iota'}\) for \(0_{\iota'}\) and the
tautologies \(\sim p\supset.\,q\supset pr\vee(\sim p)q\) and
\(\sim p\supset.\,pr\vee(\sim p)q\supset q\) in place of the two above.

\emph{(c) Clause~1.}  Let \(\mathbf H\) be
\((x_\alpha)[\mathbf F_{o\alpha}x_\alpha\equiv\mathbf G_{o\alpha}x_\alpha]\),
and write \(\mathbf G\) for \(\mathbf G_{o\alpha}x_\alpha\).  By
12\(^\alpha\) (with IV\('\) and Rule~II), \(\mathbf H\derives\mathbf F\equiv\mathbf G\).
By (a) for \(\mathbf F\) and for \(\mathbf G\), with (P), V on
\(p\equiv q\supset.\,p\supset q\),
\(\mathbf H,\mathbf F\derives\mathbf C_{\mathbf F}x_\alpha=0_{\iota'}\) and
\(\mathbf H,\mathbf F\derives\mathbf C_{\mathbf G}x_\alpha=0_{\iota'}\); then
19\(^{\iota'}\) and 20\(^{\iota'}\) (with IV\('\) and V) give
\(\mathbf H,\mathbf F\derives\mathbf C_{\mathbf F}x_\alpha=\mathbf C_{\mathbf G}x_\alpha\).
Likewise (b) gives
\(\mathbf H,\sim\mathbf F\derives\mathbf C_{\mathbf F}x_\alpha=\mathbf C_{\mathbf G}x_\alpha\).
The deduction theorem (no variable free in an assumption is substituted for
or generalized upon in these proofs) gives
\(\mathbf H\derives\mathbf F\supset\mathbf C_{\mathbf F}x_\alpha=\mathbf C_{\mathbf G}x_\alpha\)
and
\(\mathbf H\derives\sim\mathbf F\supset\mathbf C_{\mathbf F}x_\alpha=\mathbf C_{\mathbf G}x_\alpha\),
and the tautology \((p\supset r)\supset.\,(\sim p\supset r)\supset r\) gives
\(\mathbf H\derives\mathbf C_{\mathbf F}x_\alpha=\mathbf C_{\mathbf G}x_\alpha\).
The variable \(x_\alpha\) is not free in \(\mathbf H\), so VI\('\) gives
\(\mathbf H\derives(x_\alpha)[\mathbf C_{\mathbf F}x_\alpha=\mathbf C_{\mathbf G}x_\alpha]\).
In 10\(^{\iota'\alpha}\),
\((x_\alpha)[f_{\iota'\alpha}x_\alpha=g_{\iota'\alpha}x_\alpha]\supset f_{\iota'\alpha}=g_{\iota'\alpha}\),
substitute \(\mathbf C_{\mathbf F}\) and \(\mathbf C_{\mathbf G}\) for
\(f_{\iota'\alpha}\) and \(g_{\iota'\alpha}\) by IV\('\); Rule~V gives
\(\mathbf H\derives\mathbf C_{\mathbf F}=\mathbf C_{\mathbf G}\), and the
deduction theorem gives clause~1.  The axioms used are 1--6\(^\gamma\),
9\(^{\iota'}\) and 10\(^{\iota'\alpha}\).

\emph{(d) \(\derives0_{\iota'}\ne1_{\iota'}\) from Axioms 1--7.}  By
28\(^\iota\) (\cref{thm:p02-w02-p02-040-church-27-29}), whose derivation uses
Axioms 1--6\(^\gamma\) and Axiom~7 as 27\(^\iota\), and IV\('\),
\(\derives S_{\iota'\iota'}0_{\iota'}\ne0_{\iota'}\).  By Rule~II,
\(S_{\iota'\iota'}0_{\iota'}\) converts to
\(\lambda f_{\iota\iota}\lambda x_\iota(f_{\iota\iota}(0_{\iota'}f_{\iota\iota}x_\iota))\)
and then to \(\lambda f_{\iota\iota}\lambda x_\iota(f_{\iota\iota}x_\iota)\),
which is \(1_{\iota'}\); so \(\derives1_{\iota'}\ne0_{\iota'}\), and
19\(^{\iota'}\) with (P), V gives \(\derives0_{\iota'}\ne1_{\iota'}\).

\emph{(e) Clause~2.}  By (a) and the deduction theorem,
\(\derives\mathbf F\supset\mathbf C_{\mathbf F}x_\alpha=0_{\iota'}\).  By (b),
\(\sim\mathbf F,\mathbf C_{\mathbf F}x_\alpha=0_{\iota'}\derives\mathbf C_{\mathbf F}x_\alpha=1_{\iota'}\),
and 19\(^{\iota'}\), 20\(^{\iota'}\) give
\(\sim\mathbf F,\mathbf C_{\mathbf F}x_\alpha=0_{\iota'}\derives0_{\iota'}=1_{\iota'}\).
The deduction theorem twice, (d), and the tautology
\((\sim p\supset.\,q\supset r)\supset.\,\sim r\supset.\,q\supset p\) give
\(\derives\mathbf C_{\mathbf F}x_\alpha=0_{\iota'}\supset\mathbf F\).  With
\(p\supset q\supset.\,(q\supset p)\supset.\,p\equiv q\) this yields
\(\derives\mathbf C_{\mathbf F}x_\alpha=0_{\iota'}\equiv\mathbf F\), the
first formula of clause~2.

For the second, assume \(\mathbf C_{\mathbf F}=\mathbf C_{\mathbf G}\).  In
18\(^{\iota'(\iota'\alpha)}\),
\(x_{\iota'\alpha}=y_{\iota'\alpha}\supset f_{\iota'(\iota'\alpha)}x_{\iota'\alpha}=f_{\iota'(\iota'\alpha)}y_{\iota'\alpha}\),
substitute \(\mathbf C_{\mathbf F}\), \(\mathbf C_{\mathbf G}\) and
\(\lambda k_{\iota'\alpha}(k_{\iota'\alpha}x_\alpha)\) by IV\('\), with
\(k_{\iota'\alpha}\) a new variable, and convert by Rule~II: this gives
\(\mathbf C_{\mathbf F}=\mathbf C_{\mathbf G}\derives\mathbf C_{\mathbf F}x_\alpha=\mathbf C_{\mathbf G}x_\alpha\).
By the first formula of clause~2 for \(\mathbf F\) and for \(\mathbf G\),
together with 19\(^{\iota'}\), 20\(^{\iota'}\) and (P),
\(\mathbf C_{\mathbf F}=\mathbf C_{\mathbf G}\derives\mathbf F\equiv\mathbf G\):
if \(\mathbf F\) then \(\mathbf C_{\mathbf F}x_\alpha=0_{\iota'}\), hence
\(\mathbf C_{\mathbf G}x_\alpha=0_{\iota'}\), hence \(\mathbf G\), and
symmetrically.  The variable \(x_\alpha\) is not free in the assumption, so
VI\('\) and the deduction theorem give the second formula.  The axioms used
are 1--6\(^\gamma\), 7 and 9\(^{\iota'}\); Axiom~10 is not used in (d) or
(e).
\end{proof}
\end{annexitem}

\begin{annexitem}{annex:p02-w02c-p02-028-no-ax7}{thm:p02-w02c-p02-028-no-ax7}{Theorem 26 at higher types without Axiom~7}
\begin{proof}
Fix \(k\ge1\) and let \(\Sigma\) be the set of Axioms 1--6\(^\gamma\),
Axiom~8, and 10\(^{\alpha_j\alpha_j}\), 10\(^{(\alpha_j\alpha_j)(\alpha_j\alpha_j)}\)
for \(j\le k\).  Write \(\mathbf C\) for 26\(^{\alpha_k}\), a formula about
the numbers of type \(\alpha_k'=\alpha_{k+1}\), and \(\mathbf E\) for
27\(^\iota\), which is literally Axiom~7; \(\mathbf E\) has no free
variable.

\emph{Case A: from \(\mathbf E\) to \(\mathbf C\).}  By
\cref{annex:p02-w02-p02-040-church-26-30}, \(\mathbf C\) has a proof whose
axioms are variants of Axioms 1--6\(^\gamma\), 7 and 8, and Axiom~7 enters
that proof only as 27\(^\iota\).  In such a proof replace each occurrence of
a variant of Axiom~7 by the closed assumption \(\mathbf E\) followed by the
applications of Rule~I that turn \(\mathbf E\) into that variant.  The
result is a proof of \(\mathbf C\) on the assumption \(\mathbf E\) from
\(\Sigma\): its axioms are variants of members of \(\Sigma\), and the
proviso on assumptions is met because \(\mathbf E\) has no free variable.
The deduction theorem gives \(\derives_\Sigma\mathbf E\supset\mathbf C\).

\emph{Case B: from \(\sim\mathbf E\) to \(\mathbf C\).}  For a type
\(\gamma\) let \(\mathbf H_\gamma\) be the closed formula
\((x_\gamma)(y_\gamma)\,x_\gamma=y_\gamma\).
\begin{enumerate}
 \item \(\sim\mathbf E\derives_\Sigma\mathbf H_\iota\).  Unabbreviated,
   \(\mathbf E\) is \(\sim(x_\iota)\sim\sim(y_\iota)\sim\sim[x_\iota=y_\iota]\).
   From \(\sim\mathbf E\), (P), V with \(\sim\sim p\supset p\) gives
   \((x_\iota)\sim\sim(y_\iota)\sim\sim[x_\iota=y_\iota]\); 12\(^\iota\)
   (with IV\('\) and Rule~II) gives
   \(\sim\sim(y_\iota)\sim\sim[x_\iota=y_\iota]\); (P), V gives
   \((y_\iota)\sim\sim[x_\iota=y_\iota]\); 12\(^\iota\) again gives
   \(\sim\sim[x_\iota=y_\iota]\); and (P), V gives \(x_\iota=y_\iota\).
   VI\('\) on \(y_\iota\) and then on \(x_\iota\), neither free in
   \(\sim\mathbf E\), gives \(\mathbf H_\iota\).
 \item \(\mathbf H_\alpha\derives_\Sigma\mathbf H_{\alpha\beta}\) whenever
   10\(^{\alpha\beta}\in\Sigma\).  From \(\mathbf H_\alpha\), two uses of
   12\(^\alpha\), substituting \(f_{\alpha\beta}x_\beta\) and
   \(g_{\alpha\beta}x_\beta\) by IV\('\), give
   \(f_{\alpha\beta}x_\beta=g_{\alpha\beta}x_\beta\).  VI\('\) on
   \(x_\beta\) gives \((x_\beta)[f_{\alpha\beta}x_\beta=g_{\alpha\beta}x_\beta]\),
   and 10\(^{\alpha\beta}\) with Rule~V gives
   \(f_{\alpha\beta}=g_{\alpha\beta}\).  VI\('\) on \(g_{\alpha\beta}\) and
   then on \(f_{\alpha\beta}\) gives \(\mathbf H_{\alpha\beta}\).  No
   variable generalized upon is free in the closed assumption.
 \item Hence \(\sim\mathbf E\derives_\Sigma\mathbf H_{\alpha_{j}}\) for
   \(j\le k+1\), by induction on \(j\): the case \(j=0\) is item~1, and since
   \(\alpha_{j+1}=((\alpha_j\alpha_j)(\alpha_j\alpha_j))\), item~2 applied
   with 10\(^{\alpha_j\alpha_j}\) gives \(\mathbf H_{\alpha_j\alpha_j}\) and,
   applied again with 10\(^{(\alpha_j\alpha_j)(\alpha_j\alpha_j)}\), gives
   \(\mathbf H_{\alpha_{j+1}}\).  The assumption is closed, so proofs on it
   may be chained.
 \item From \(\mathbf H_{\alpha_{k+1}}\), two uses of 12\(^{\alpha_{k+1}}\)
   give \(x_{\alpha_{k+1}}=y_{\alpha_{k+1}}\), and (P), V with the tautology
   \(p\supset.\,q\supset.\,r\supset.\,s\supset p\) gives
   \[
   \begin{aligned}
    &N_{o\alpha_{k+1}}x_{\alpha_{k+1}}\supset.\,N_{o\alpha_{k+1}}y_{\alpha_{k+1}}\supset.\\
    &\qquad S_{\alpha_{k+1}\alpha_{k+1}}x_{\alpha_{k+1}}=S_{\alpha_{k+1}\alpha_{k+1}}y_{\alpha_{k+1}}
    \supset x_{\alpha_{k+1}}=y_{\alpha_{k+1}},
   \end{aligned}
   \]
   which is \(\mathbf C\).
\end{enumerate}
So \(\sim\mathbf E\derives_\Sigma\mathbf C\), and the deduction theorem
gives \(\derives_\Sigma\sim\mathbf E\supset\mathbf C\).

\emph{Conclusion.}  The tautology
\((p\supset r)\supset.\,(\sim p\supset r)\supset r\), with (P) and two
applications of Rule~V to the results of Cases A and B, gives
\(\derives_\Sigma\mathbf C\).  Axiom~7 is not in \(\Sigma\).
\end{proof}
\end{annexitem}
